%% sections/04-evaluation.tex

\section{Evaluation}
\label{sec:evaluation}

We evaluate the innovation-cluster-driven amendment pipeline on three properties:
(1) the self-calibration claim (innovation rate decays as the rubric matures and
reactivates on novel archetypes); (2) cluster coherence (amendment-triggering
clusters are thematically unified); and (3) dimension distribution (the Atmospheric
Landing dimension requires disproportionately many amendments, for structural
rather than incidental reasons). We report against a 49-session, 7-campaign corpus
spanning 4 complete days of production use.

\subsection{Innovation Rate by Session and Campaign}

Table~\ref{tab:innovations} reports innovations logged per campaign across the
7-campaign corpus, alongside rubric version at campaign open and close.

\begin{table}[t]
\caption{Innovations Logged Per Campaign: Session Range, Total, Mean per Session}
\label{tab:innovations}
\begin{tabular}{@{}lllrrl@{}}
\toprule
Campaign & Rubric open & Rubric close & Total & Mean/session & Archetype \\
\midrule
C1 & v1.0 & v1.4 & 40+ & 5.7 & grief-themed \\
C2 & v1.4 & v1.6 & 12  & 1.7 & covenant-restoration \\
C3 & v1.6 & v1.8 & 16  & 2.3 & faith/authority \\
C4 & v1.8 & v1.12 & 16 & 2.3 & siege-resource \\
C5 & v1.12 & v1.12 & 8  & 1.1 & professional-code \\
C6 & v1.12 & v1.12 & 15 & 2.1 & duty-spine \\
C7 & v1.12 & v1.12 & 7  & 1.0 & no-designed-stakes \\
\midrule
Total & & & 114+ & 2.3 & \\
\bottomrule
\end{tabular}
\end{table}

The self-calibration pattern is clearly visible. C1 has the highest mean
innovation rate (5.7 per session) because the rubric begins at v1.0 with no
accumulated anchors, and every novel quality phenomenon produces a log entry.
C2 drops sharply to 1.7 per session: the four C1 amendments absorbed
per-PC reception tracking, cross-character compound mechanics, anchor-level
fallback requirements, and finder/receiver divergence --- the highest-frequency
gaps in a v1.0 rubric applied to grief-themed TRPG sessions. C3 and C4 produce
moderate bursts (2.3 per session each) as their archetypes introduce quality
phenomena the v1.4 and v1.6 anchors had not seen. C5 through C7, all running
under the mature v1.12 rubric, produce the lowest per-session rates (1.0--2.1).

Figure~\ref{fig:session-rate} shows innovation counts per session for C1
(high-innovation baseline), illustrating the within-campaign decay from
S01 through S07.

\begin{figure}[t]
  \centering
  \fbox{\parbox{0.85\columnwidth}{\centering
    [Bar chart: innovations per session, C1 (S01--S07). \\
    S01: 12. S02: 8. S03: 5. S04: 5. S05: 4. S06: 3. S07: 3. \\
    Total: 40. Mean: 5.7. \\
    Arrow annotation: ``v1.1 adopted after S01. v1.2 after S02. \\
    v1.3 after S03. v1.4 after S04.'']}}
  \caption{Innovations per session in C1 (grief-themed, 7 sessions).
    The within-campaign decay tracks the amendment adoption events:
    each amendment absorbs a cluster of recurring patterns, reducing
    the gap between rubric and observed outcomes.}
  \label{fig:session-rate}
\end{figure}

Notable: C6 (military-unit, duty-spine) produces 15 innovations despite running
under the mature v1.12 rubric. Inspection of the C6 log reveals that the
Unit Cohesion Die mechanic produced quality phenomena absent from all prior
campaigns: ``grief-paragraph resource expression'' (resource decisions shaped
by grief-paragraph commitment rather than tactical calculation),
``unit-member-as-named-loss'' atmospheric landing, and ``order-defiance as
character commitment.'' All three were classified as universal-scope innovations
--- not party-specific or adventure-specific --- but they required the
combination of duty-spine archetype and Cohesion Die mechanic to appear. Neither
element alone had appeared in the corpus before C6.

\subsection{Amendment Adoption Rate and Archetype Reactivation}

Figure~\ref{fig:amendment-rate} shows the amendment adoption rate (amendments
adopted per session) by campaign block.

\begin{figure}[t]
  \centering
  \fbox{\parbox{0.85\columnwidth}{\centering
    [Line chart: amendments adopted per session window. \\
    C1 (S01--S07): 0.57 per session (4 adopted). \\
    C2 (C2-S01 to C2-S07): 0.29 per session (2 adopted). \\
    C3 (C3-S1 to C3-S7): 0.29 per session (2 adopted). \\
    C4 (C4-S1 to C4-S7): 0.57 per session (4 adopted). \\
    C5--C7 (21 sessions combined): 0.00 per session (0 adopted). \\
    Dashed line: 0.25 per session (approximate early steady-state).]}}
  \caption{Amendment adoption rate by session position.
    C1 and C4 share the highest rate (0.57/session). C2 and C3 stabilize
    at 0.29/session. C5--C7 approach zero as v1.12 matures.}
  \label{fig:amendment-rate}
\end{figure}

The decay-and-reactivation pattern confirms Claim 1. C1's rate (0.57/session)
reflects a fresh rubric encountering 7 sessions of quality outcomes for the first
time. C2 drops to 0.29 as the rubric absorbs the most common patterns. C3 holds
at 0.29 as the faith/authority archetype introduces 2 new clusters
(framework-failure arc-activation from a fumbled Insight check; behavioral-stance
arc-activation from combat-as-precondition). C4 reactivates at 0.57 as the
siege-resource mechanic produces 4 new clusters the grief-adjacent and covenant
anchors had not addressed.

C5 through C7, all under v1.12, produce no adoptions. This does not imply the
rubric is complete for all conceivable archetypes; it implies that the three
archetypes run under v1.12 produced innovations classified as either
adventure-specific or party-specific in scope, or as already captured by
existing v1.12 anchors. No universal-scope innovation from C5--C7 went
unadopted.

\subsection{Cluster Coherence}

Table~\ref{tab:cluster-coherence} characterizes all 15 amendment-triggering
dimension-clusters by cluster size, session span, and thematic label.

\begin{table}[t]
\caption{Cluster Coherence: All 15 Amendment-Triggering Dimension-Clusters}
\label{tab:cluster-coherence}
\begin{tabular}{@{}llrcl@{}}
\toprule
Version & Dim & Size & Sessions & Thematic label \\
\midrule
v1.1 & MF  & 2 & S01       & Cross-dimensional harmlessness \\
v1.1 & AL  & 3 & S01       & Per-PC reception + chain-reaction \\
v1.2 & MFid & 4 & S01--S02 & Anchor-level single-gate failure \\
v1.2 & AL  & 3 & S02       & Finder/receiver divergence + silent reception \\
v1.3 & MFid & 3 & S03      & Positional-geometry failure \\
v1.4 & MF  & 2 & S03--S04  & Cross-character compound \\
v1.4 & AL  & 2 & S04--S05  & Multi-scene silent-reception chain \\
v1.5 & AL  & 3 & S07, C2-S01--S02 & Cross-session act-without-announcement \\
v1.6 & AL  & 2 & S07, C2-S04 & Simultaneous arc-convergence \\
v1.7 & AL  & 2 & C3-S01    & Framework-failure arc-activation \\
v1.8 & AL  & 2 & C3-S03--S04 & Behavioral-stance arc-activation \\
v1.9 & AL, CA & 4 & C4-S01 & Grief-paragraph resonance + sustained pattern \\
v1.10 & CA & 2 & C4-S02--S03 & Character-register conversion \\
v1.11 & MF & 2 & C4-S03--S04 & Portent-as-directed-foreknowledge \\
v1.12 & AL & 2 & C4-S02--S03 & Intra-session restraint \\
\bottomrule
\end{tabular}
\end{table}

All 15 clusters exhibited thematic unity: innovations within a cluster addressed
the same underlying quality phenomenon from different session contexts. No cluster
contained innovations addressing genuinely distinct phenomena that happened to
share a dimension tag by accident.

This result supports Claim 2. The coherence is particularly striking in the
Atmospheric Landing clusters, where 9 clusters (v1.1, v1.2, v1.4, v1.5, v1.6,
v1.7, v1.8, v1.9, v1.12) each addressed clearly distinguishable phenomena:
who receives a beat (v1.1), who finds versus who receives it (v1.2), how
reception persists across scenes (v1.4), how reception expresses across sessions
(v1.5), how two PC arcs converge simultaneously (v1.6), how framework failure
delivers atmospheric content (v1.7), how behavioral stance activates an arc
(v1.8), how grief shapes resource decisions (v1.9), and how deferred expression
within a session scores differently from immediate reception (v1.12). These are
nine distinct quality phenomena, all in the same dimension, each the product of
a separate innovation cluster.

\subsection{Innovation Scope Distribution}

Figure~\ref{fig:taxonomy} shows the scope distribution across 120 analyzed
innovation entries (120+ total; 120 with scope classification).

\begin{figure}[t]
  \centering
  \fbox{\parbox{0.85\columnwidth}{\centering
    [Bar chart: innovations by scope classification. \\
    Universal: 78 (65\%). \\
    Party-specific: 24 (20\%). \\
    Adventure-specific: 18 (15\%). \\
    Total analyzed: 120. \\
    All 15 amendment clusters contained $\geq$1 universal-scope innovation.]}}
  \caption{Innovation scope distribution. Universal-scope innovations
    are the primary raw material for amendments; party-specific and
    adventure-specific innovations provide corroborating evidence only.}
  \label{fig:taxonomy}
\end{figure}

Of 120 analyzed innovations, 78 (65\%) were classified as Universal. All 15
amendment-triggering clusters contained at least one Universal-scope innovation;
no amendment was adopted on party-specific or adventure-specific evidence alone.
This validates the scope field as a necessary gate: 24 party-specific and 18
adventure-specific innovations were logged without triggering amendments, which
is the intended behavior (party-specific patterns do not generalize; they belong
in the party documentation, not the rubric).

\subsection{Dimension Distribution: The Atmospheric Landing Asymmetry}

Table~\ref{tab:dimension-distribution} reports amendment counts by dimension
across all 13 versions.

\begin{table}[t]
\caption{Amendment Count by Dimension}
\label{tab:dimension-distribution}
\begin{tabular}{@{}lrc@{}}
\toprule
Dimension & Amendments & Versions \\
\midrule
Atmospheric Landing (AL)  & 9 & v1.1--v1.12 \\
Mechanical Fairness (MF)  & 3 & v1.1, v1.4, v1.11 \\
Module Fidelity (MFid)    & 2 & v1.2, v1.3 \\
Character Agency (CA)     & 2 & v1.9, v1.10 \\
Engagement (EG)           & 0 & --- \\
Pacing (PA)               & 0 & --- \\
Surprise (SU)             & 0 & --- \\
Table Readiness (TR)      & 0 & --- \\
\midrule
Total                     & 16 & (15 clusters; v1.9 dual) \\
\bottomrule
\end{tabular}
\end{table}

Atmospheric Landing received 9 of the 16 dimension-amendments across 13 version
increments. This asymmetry has a structural explanation. Engagement, Pacing,
Surprise, and Table Readiness all score bounded properties: either PCs drove
scenes or they did not; either pacing transitions were clean or they were not;
either 7+ surprise tags fired or they did not; either the DM required outside
reference or they did not. Initial anchors for these dimensions are adequate
because the quality space is effectively finite.

Mechanical Fairness and Module Fidelity required 3 and 2 amendments respectively,
primarily because their 10-band definitions (the most sophisticated anchor band)
required progressive refinement as the corpus revealed edge cases. Character Agency
required 2 amendments because the initial anchors assumed agency would manifest
as explicit choices; the corpus revealed that sustained behavioral patterns and
character-register conversion are equally valid agency expressions.

Atmospheric Landing required 9 amendments because it spans the full space of
how emotional and narrative intent can be realized in session play. This space is
effectively unbounded: every campaign archetype that introduces a new relational
mechanic or grief-configuration produces at least one quality phenomenon not
captured by existing anchors. The initial AL anchors were necessary foundations
(reception tracking, chain-reaction scoring), but each campaign archetype added
a new realization mode.

\subsection{C5--C7 Innovation Bursts Confirm Archetype-Independence}

Campaigns 5 through 7 were designed specifically to test the rubric against novel
archetypes: professional-code (duty without explicit grief-hooks), duty-spine
(military-unit relational mechanics), and no-designed-stakes (player-driven content
creation without designed emotional anchors). All three ran under the stable v1.12
rubric, which had absorbed C1--C4's innovations.

C5 (professional-code) produced 8 innovations, all classified as either
party-specific or adventure-specific. This is the expected result for an archetype
whose quality patterns are structurally similar to C1--C4's grief-adjacent patterns
but party-specific in their expression.

C6 (duty-spine, Unit Cohesion Die) produced 15 innovations --- the highest count
of any campaign after C1 --- and all 15 were classified as universal-scope. Despite
not triggering amendments (because the v1.12 anchors were found to cover the
patterns on closer review), these innovations confirm that novel mechanics with
novel relational consequences reliably produce innovation bursts. The C6 burst
is evidence that the pipeline remains sensitive under a mature rubric: it did
not miss the novelty; it identified it, classified it, and correctly determined
that existing anchors were adequate after all.

C7 (no-designed-stakes) produced 7 innovations, confirming the floor: even a
campaign with no designed emotional content produces a small number of universal
innovations as the AI-simulated players invent their own emotional stakes.

Taken together, C5--C7 validate a prediction of the mechanism: a mature rubric
does not become insensitive; it simply stops amending because the patterns it
encounters are either already covered or not universal. The innovation rate
stabilizes at 1--2 per session as a floor, not zero, indicating that some
portion of creative session outcomes will always exceed any finite rubric.
